\chapter{Conclusion and Future Work}
\label{ch:conclusion}

% ===========================================================================
% WARNING. The original skeleton claimed contribution 1 as "EMPIRICAL:
% LLM-agent coordination degrades more gracefully than a doctrine-grounded
% workflow baseline under novel disruption." THAT IS THE OPPOSITE OF THE
% RESULT. It was the proposal's expectation. It has been replaced below.
% ===========================================================================

\section{Summary}
\label{sec:concl:summary}
% Restate the four research questions and answer each in two or three
% sentences. No new evidence here, and no hedging.
%   RQ1 Can it beat the scripted baseline? No, and it also loses to a null
%       policy. Chapters~\ref{ch:results} and \ref{ch:critique}.
%   RQ2 Why not? A commitment failure, not a comprehension failure.
%       Chapter~\ref{ch:discussion}.
%   RQ3 What would have to change? Constrain commitment; test on scenarios
%       that discriminate. Section~\ref{sec:concl:future}.
%   RQ4 What should a designer consider? Chapter~\ref{ch:design}.
% Then the four engineering sub-questions, one line each, pointing at
% Chapters~\ref{ch:data}, \ref{ch:agentic}, \ref{ch:eval} and \ref{ch:results}.

\section{Contributions Revisited}
\label{sec:concl:contributions}
% The four contributions from Section~\ref{sec:intro:contributions}, restated
% against the evidence that was actually produced:
%
%  1. ARTEFACT. Built and running: 10 agents across 8 organisations, Postgres
%     as single source of truth and active broker, plus a doctrine-grounded
%     baseline at audited parity (Chapter~\ref{ch:validity}).
%
%  2. EMPIRICAL, AND NEGATIVE. 310 scored runs over 7 scenarios. The agentic
%     architecture does not outperform the scripted baseline on resident
%     survival, and underperforms a null policy. This is a result, and it is
%     reported as one.
%
%  3. DIAGNOSTIC. The deficit is located and measured: commitment, not
%     comprehension. This is the part a future builder can act on.
%
%  4. METHODOLOGICAL. The shared append-only event log as the backbone of a
%     fair, auditable, replayable comparison; and Gate 0, the treatment-
%     validity check that catches a non-discriminating disruption before it is
%     scored. A protocol that found its own primary arm degenerate and
%     specified the gate that would have caught it.

\section{Future Work}
\label{sec:concl:future}
% RQ3 lives here. Lead with the one experiment that matters and rank the rest.
%
% 1. THE COMMITMENT-CONSTRAINED VARIANT. The highest-value follow-up this
%    analysis produces. Cap re-plans, or forbid cancelling a leg already in
%    flight; run against the same scenarios under the same models. One
%    campaign, no new scenario authoring. It is also the experiment that would
%    settle the causal direction this thesis explicitly leaves open.
%
% 2. SCENARIOS THAT PASS GATE 0. Author disruptions where a no-replan control
%    demonstrably fails, so adaptation has something to prove. Until that runs,
%    "can an agentic architecture ever beat a scripted one" is open, not
%    answered.
%
% 3. A DESIGN THE EVALUATION COULD SEE. A larger k (10 could not detect a
%    moderate effect), more than one SIM_TIME_SCALE (H4 was never tested), and
%    a model set fixed before the arm runs.
%
% 4. Longer-horizon: resident-level agents, an equity dimension in the rubric,
%    compound and cascading disruptions, multi-case replication beyond a single
%    facility, other cities and hazards.
